\scanpage{120}
$q\mid D$ non nul fibre par fibre.
% slanted
\margin{vérifier cette affirmation !}
ou encore, en identifiant $q\mid D$ à un hom
\[
q^{\unsure{\prime}}:D^{\otimes2}\longrightarrow L
\]
que cet hom est un \uline{iso}
\[
q:D^{\otimes2}\xrightarrow{\sim}L\tag*{(6)}
\]
\deleted{La condition} En combinant avec l'iso. (5), on trouve un iso.
\[
L^{\otimes2}\simeq L
\]
i.e. un iso
\[
D\simeq L\simeq\mathcal O_S\tag*{(7)}
\]
Ce qu'on peut aussi exprimer ainsi :

\uline{Proposition} \inserted{Pour que $H=\mathbb P(V)$ soit non tangent à $Q\subset P$, il f et s} qu'\unsure{il|Il} existe une \deleted{unique} section $a$ de $T$ \deleted{$\unsure{V^\vee|V}$}, et une section $q$ du sous-fibré vectoriel $L^\vee$ de $\operatorname{Sym}^2(\check T)$ \unsure{définiss\textsuperscript{t}|définissant} $Q$, tels que

% margin note written sideways
\margin{Ceci équivaut aussi à : $q\mid\unsure{V^\perp|V}$ \uline{\unsure{régul.|régulière}}}
\[
\text{a)}\quad q(a)=1\tag*{(8)}
\]
\[
\text{b)}\quad\deleted{\unsure{E=\{|E=}}\pi(x)\unsure{=|\in}\deleted{\unsure{\Gamma(-,T)\mid|\Gamma(-,T)}}\varphi(a,x)\deleted{\unsure{=1\}}}\quad(x\in\Gamma(-,T)\unsure{)}
\tag*{(9)}
\]
\inserted{où $\varphi$ est la forme bil. associée à $q$.}

Alors $a,q$ sont déterminés de façon unique, et ce sont respectivement des bases de $D$ et de $L^\vee$. Enfin, on a
\[
\pi(a)=2\tag*{(10)}
\]
Il reste à prouver cette dernière relation. \deleted{Or \unsure{on a} plus gén, pour $x\in\Gamma(\unsure{-,|,}T)$}
% archive sheet number: 119
